% ============================================================
%  APPENDIX TO CHAPTER 3 -- EXERCISES S3.1-S3.18
%  Imported from Tao I, ECNU, USTC, Yu Pin.  Not Rudin's.
% ============================================================
\ssection{Appendix exercises}

\begin{roadmap}[How these differ from Rudin's twenty-five]
Rudin's own Exercises 3.1--3.25 end the chapter proper. These eighteen come
from the four comparison texts.

\medskip
\raggedright
\textbf{Group I (1--5): sequences.} Coupled recursions --- Gauss's
arithmetic--geometric mean among them --- where Rudin's 3.16--3.17 use a
single recursion; plus the means package and bounded variation.

\medskip
\textbf{Group II (6--12): series.} The decay rate of the terms of a
convergent series, Cauchy--Schwarz in $\ell^2$, infinite products, and the
positive/negative split that Rudin buries inside the proof of 3.54.

\medskip
\textbf{Group III (13--16): upper and lower limits.} The algebra of
Supplement A, put to work.

\medskip
\textbf{Group IV (17--18): summation over sets.} Iterated sums, and the
counterexample that shows absolute convergence is not decoration.

\medskip
\textbf{Marked $\star$:} harder, or requiring a result from a later chapter.
\end{roadmap}

\begin{enumerate}[label=\textbf{S3.\arabic*.}, leftmargin=3.4em, itemsep=9pt,
                  topsep=8pt, labelsep=0.5em]

% ---------------- Group I ----------------
\item \emph{(Arithmetic--geometric mean.)} Let $a_1 > b_1 > 0$ and put
\[ a_{n+1} = \frac{a_n+b_n}{2}, \qquad b_{n+1} = \sqrt{a_nb_n}. \]
Prove that $\lim a_n$ and $\lim b_n$ both exist and are
equal.\kn{Gauss's AGM. The engine is the inequality $\sqrt{xy} \le (x+y)/2$,
which gives $b_n < b_{n+1} < a_{n+1} < a_n$: the two sequences are monotone
in opposite directions and interleaved, so 3.14 applies to each and the gap
$a_{n+1}-b_{n+1} \le (a_n-b_n)/2$ closes. Rudin's Exercises 3.16--3.17 are
single recursions converging to a square root; a \emph{coupled} recursion is
materially different, and the common limit here has no elementary closed
form --- it is an elliptic integral.}
\provenance{ECNU \cjk{习题}{Ex} 2.3 \#11 (bk.\ p.\,38).}

\item \emph{(Arithmetic--harmonic mean.)} With $a_1, b_1 > 0$ put
$a_{n+1} = (a_n+b_n)/2$ and $b_{n+1} = 2a_nb_n/(a_n+b_n)$. Prove both limits
exist and equal $\sqrt{a_1b_1}$.\kn{The companion to exercise 1, and this
time the answer \emph{is} elementary --- because $a_nb_n$ is invariant:
$a_{n+1}b_{n+1} = a_nb_n$ for every $n$. Spot the invariant and the problem
collapses. Compare the strategy index's second pattern in Chapter 1: hunt for
the quantity the recursion does not move.}
\provenance{ECNU \cjk{总练习题}{review} 2 \#8 (bk.\ p.\,39).}

\item Prove that $a_n \to a$ implies $(a_1+\cdots+a_n)/n \to a$, and that if
also $a_n > 0$ then $\sqrt[n]{a_1\cdots a_n} \to a$. Deduce that
$b_{n+1}/b_n \to a$ implies $\sqrt[n]{b_n} \to a$ for $b_n>0$, and use it to
evaluate $\lim n/\sqrt[n]{n!}$.\kn{Rudin gives only the first part, as
Exercise 3.14(a). The deduction is the content of 3.37 --- the root test
dominates the ratio test --- and the last limit is $e$, obtained by applying
(c) to $b_n = n^n/n!$. Show also that the converses all fail: averaging
destroys information, so $(-1)^n$ has Ces\`aro mean $0$ without converging.}
\provenance{ECNU \cjk{总练习题}{review} 2 \#3--\#4 (bk.\ p.\,39); USTC
Ch.\ 1 suppl.\ \#6--\#11.}

\item Suppose $\sum_{k=1}^{n} |a_{k+1}-a_k| \le M$ for every $n$. Prove that
$(a_n)$ converges.\kn{Bounded variation implies convergence --- the sequence
version of a notion Chapter 6 will develop for functions. The proof is one
line once you see it: $a_n = a_1 + \sum_{k<n}(a_{k+1}-a_k)$ exhibits $(a_n)$
as the partial sums of an absolutely convergent series. Note the converse
fails: a convergent sequence can wander infinitely far, as
$a_n = \sum_{k\le n} (-1)^k/k$ shows.}
\provenance{ECNU \cjk{总练习题}{review} 2 \#6 (bk.\ p.\,39); USTC Ch.\ 1
suppl.\ \#5.}

\item Prove that if every convergent subsequence of a bounded sequence has
the same limit $L$, then the sequence itself converges to $L$. Deduce that a
bounded divergent sequence has two subsequences with different
limits.\kn{The cleanest packaging of Rudin's 3.6--3.7. Argue by
contradiction: if $x_n \not\to L$, infinitely many terms stay $\varepsilon$
away, and that subsequence is bounded, so by 3.6(b) it has a convergent
sub-subsequence --- whose limit is not $L$. Note the hypothesis
\emph{bounded} is essential: $x_n = n$ has no convergent subsequence at all,
so the condition holds vacuously.}
\provenance{USTC \S1.2 (bk.\ p.\,15), left there as an exercise.}

% ---------------- Group II ----------------
\item Prove that if $\Sigma u_n$ is a convergent series of positive terms and
$(u_n)$ is monotone, then $\lim n\,u_n = 0$.\kn{Sharper than 3.23, which
gives only $u_n \to 0$. The trick is to look at a block: monotonicity makes
$u_{n+1}+\cdots+u_{2n} \ge n\,u_{2n}$, and the left side tends to $0$ by the
Cauchy criterion. Note what fails without monotonicity --- one can have a
convergent series with $n u_n \not\to 0$ by putting spikes at sparse indices.
This is the quantitative reason the harmonic series is the boundary.}
\provenance{ECNU \cjk{总练习题}{review} 12 \#1 (\cjk{下}{vol.2}
bk.\ p.\,24).}

\item Prove that if $\Sigma a_n^2$ and $\Sigma b_n^2$ converge then
$\Sigma a_nb_n$ and $\Sigma(a_n+b_n)^2$ converge, and
\[ \Bigl(\sum a_nb_n\Bigr)^2 \le \sum a_n^2 \sum b_n^2, \]
\[ \Bigl(\sum (a_n+b_n)^2\Bigr)^{1/2} \le
   \Bigl(\sum a_n^2\Bigr)^{1/2} + \Bigl(\sum b_n^2\Bigr)^{1/2}. \]\kn{%
Cauchy--Schwarz and Minkowski in $\ell^2$ --- Theorem 1.35 and 1.37(e) in
infinite dimensions, which Rudin never revisits after Chapter 1. Prove the
finite versions and pass to the limit; the only care needed is that the
partial sums are increasing and bounded. This is the first appearance of the
space $\ell^2$, which Chapter 11 treats seriously.}
\provenance{ECNU \cjk{总练习题}{review} 12 \#8 (\cjk{下}{vol.2}
bk.\ p.\,24).}

\item Let $a_n > 0$. Prove that $\prod_{k=1}^n (1+a_k)$ converges (to a
finite nonzero limit) if and only if $\Sigma a_n$ converges. Evaluate
$\prod_{k=0}^{\infty}\bigl(1+q^{2^k}\bigr)$ for $|q|<1$.\kn{Infinite
products reduced entirely to Chapter 3 by taking logarithms: convergence of
$\Sigma \log(1+a_k)$ is equivalent to convergence of $\Sigma a_k$, since
$\log(1+x)\sim x$. The evaluation telescopes --- multiply by $(1-q)$ and
watch $(1-q)(1+q)(1+q^2)(1+q^4)\cdots = 1-q^{2^{n}}$ --- giving $1/(1-q)$.
None of the four texts devotes a section to infinite products; this exercise
is the whole theory in miniature.}
\provenance{ECNU \cjk{习题}{Ex} 12.2 \#15 (\cjk{下}{vol.2} bk.\ p.\,17);
USTC \cjk{习题}{Ex} 1.1 \#8(5).}

\item Let $a_n \ge 0$ with $(na_n)$ bounded. Prove $\Sigma a_n^2$ converges.
Prove also that if $\Sigma u_n$ is a convergent positive series then
$\Sigma\sqrt{u_nu_{n+1}}$ converges.\kn{Two decay-rate arguments. The first
is comparison with $\Sigma 1/n^2$; the second is the arithmetic--geometric
mean inequality, $\sqrt{u_nu_{n+1}} \le (u_n+u_{n+1})/2$. Both illustrate the
same habit: when a series resists direct attack, bound its terms by
something already known to converge, and prefer the crudest bound that
works.}
\provenance{ECNU \cjk{习题}{Ex} 12.2 \#5, \#7 (\cjk{下}{vol.2}
bk.\ p.\,16).}

\item With $p_n = (|u_n|+u_n)/2$ and $q_n = (|u_n|-u_n)/2$, prove that if
$\Sigma u_n$ converges conditionally then $\Sigma p_n$ and $\Sigma q_n$ both
diverge.\kn{The lemma inside Rudin's proof of 3.54, isolated. If both
converged, $\Sigma|u_n| = \Sigma(p_n+q_n)$ would converge, contradicting
conditionality; if exactly one converged, $\Sigma u_n = \Sigma(p_n-q_n)$
would diverge. So both piles are infinite --- which is precisely the fuel
Riemann's rearrangement burns. Worth extracting because the statement, not
the rearrangement, is what makes 3.54 believable.}
\provenance{ECNU \cjk{习题}{Ex} 12.3 \#4 (\cjk{下}{vol.2} bk.\ p.\,23).}

\item $\star$~Prove that $\displaystyle\sum_{n=1}^{\infty}
\frac{(-1)^{\lfloor\sqrt n\rfloor}}{n}$ converges.\kn{The signs come in runs
of lengths $3,5,7,\dots$, so the alternating series test (3.43) does not
apply at all. Group the terms into blocks of constant sign, show the grouped
series converges by comparison, and then use the fact that the terms tend to
$0$ to remove the parentheses --- which is legitimate here precisely because
the block lengths grow slowly. A genuine stress test of S3.5 and S3.8.}
\provenance{ECNU \cjk{习题}{Ex} 12.3 \#8 (\cjk{下}{vol.2} bk.\ p.\,24).}

\item Given a convergent series of positive terms $\Sigma u_n$, must there
exist $\varepsilon > 0$ with $u_n/n^{-(1+\varepsilon)} \to c$ for some
$c > 0$?\kn{No --- and the exercise is the concrete form of ``there is no
slowest convergent series,'' the moral of Supplement B. Convergent series can
decay arbitrarily slowly, so no single power of $n$ serves as a universal
gauge. This is why the Raabe/Gauss hierarchy has no last member, and why
Rudin's remark after 3.35 is a structural fact rather than an omission.}
\provenance{ECNU \cjk{总练习题}{review} 12 \#5(3) (\cjk{下}{vol.2}
bk.\ p.\,24).}

% ---------------- Group III ----------------
\item Prove the parts of S3.3: $\liminf a_n = -\limsup(-a_n)$;
$\liminf a_n + \liminf b_n \le \liminf(a_n+b_n)$; and, for positive
sequences with $\liminf a_n > 0$, $\limsup(1/a_n) = 1/\liminf a_n$.\kn{All
three fall out of the $\inf$-of-$\sup$s formula S3.1 plus the elementary
facts $\sup(-A) = -\inf A$ and $\inf(A+B) \ge \inf A + \inf B$. Doing them
once makes clear why the inequalities point the way they do: taking a
supremum over a tail is a \emph{sublinear} operation, and sublinearity is
one-directional.}
\provenance{ECNU \cjk{习题}{Ex} 7.2 \#2 (bk.\ p.\,159).}

\item Let $a_n > 0$ and suppose $\limsup a_n \cdot \limsup(1/a_n) = 1$.
Prove that $(a_n)$ converges.\kn{A one-line consequence of S3.3(e) once you
see it: the hypothesis says $\limsup a_n = \liminf a_n$, and 3.17 finishes
it. The pleasure is that a statement about products of upper limits secretly
encodes ``the two limits coincide.'' Note the hypothesis forces
$\liminf a_n > 0$, so the reciprocal rule applies.}
\provenance{ECNU \cjk{习题}{Ex} 7.2 \#4 (bk.\ p.\,159).}

\item $\star$~Suppose $\liminf x_n = A < B = \limsup x_n$ and
$\lim(x_{n+1}-x_n) = 0$. Prove that the set of subsequential limits of
$(x_n)$ is exactly $[A,B]$.\kn{A beautiful result with no analogue anywhere
in Rudin. The hypothesis says the sequence moves in ever-smaller steps while
still travelling between $A$ and $B$ infinitely often --- so it cannot skip
any value in between. Formally: given $c \in (A,B)$ and $\varepsilon>0$,
find a stretch of indices crossing from below $c$ to above $c$; because the
steps are eventually smaller than $\varepsilon$, some term of that stretch
lands within $\varepsilon$ of $c$. Compare S2.10 in the Chapter 2 appendix,
where consecutive-terms-close fails to give convergence: here the same
hypothesis is turned into positive information.}
\provenance{ECNU \cjk{总练习题}{review} 7 \#3 (bk.\ p.\,160), starred there
as well.}

\item Show by example that $\limsup(a_nb_n) = \limsup a_n \cdot \limsup b_n$
can fail for positive bounded sequences, and prove it holds when one of the
sequences converges.\kn{The multiplicative twin of the pitfall at S3.3.
Take $a_n$ alternating between $1$ and $2$ and $b_n$ alternating the other
way: each has $\limsup = 2$, but the product is constantly $2$. The repair
is the same --- a convergent factor peaks everywhere at once, so it cannot
dodge the other sequence's peaks.}
\provenance{ECNU \cjk{习题}{Ex} 7.2 \#2 (bk.\ p.\,159), product part.}

% ---------------- Group IV ----------------
\item Define $f : \N\times\N \to \R$ by $f(n,n) = 1$, $f(n,n+1) = -1$, and
$f = 0$ otherwise. Compute both iterated sums
$\sum_n\bigl(\sum_m f(n,m)\bigr)$ and $\sum_m\bigl(\sum_n f(n,m)\bigr)$, and
explain why S3.12 does not apply.\kn{The answers are $0$ and $1$. Every row
and every column is a \emph{finite} sum, so nothing diverges and nothing is
subtle --- and still the order of summation matters, because
$\sum_{(n,m)}|f(n,m)|$ is infinite. The lesson is that absolute convergence
is the entire hypothesis of Fubini, not a technical convenience. Rudin's
Exercise 3.25 asks for a relative of this.}
\provenance{Tao I, \S8.2 discussion (bk.\ p.\,188).}

\item $\star$~Let $X$ be any set and $f : X \to \R$ with
$\sum_{x\in X}|f(x)|$ finite. Prove that $\{x : f(x) \ne 0\}$ is at most
countable, and that $\sum_{x\in X} f(x)$ is independent of the enumeration
used to compute it.\kn{S3.11 and the well-definedness behind S3.10. For the
first part, show $\{x : |f(x)| > 1/n\}$ is finite for each $n$ and take the
union. For the second, the two enumerations differ by a bijection of $\N$,
so this is Rudin's 3.55 --- which is exactly the point: absolute convergence
is what makes ``the sum over a set'' meaningful, and 3.54 says the notion
collapses without it.}
\provenance{Tao I, Lemma 8.2.5 and Ex.\ 8.2.2 (bk.\ p.\,188).}

\end{enumerate}
